% Classic invoice - Pressa public template (MIT). https://pressa.dev/templates/invoice?design=classic-invoice
% Line totals, subtotal, tax and total are computed from the data, each rounded to the
% currency's minor unit (cents, whole yen) as printed, so the printed numbers add up.
% Unit prices keep up to 4 decimals for items priced below a cent.
\documentclass[10pt,a4paper]{article}
\usepackage[T1]{fontenc}
\usepackage{charter}
\usepackage[margin=20mm,top=18mm]{geometry}
\usepackage[table]{xcolor}
\usepackage{tabularx}
\usepackage{longtable}
\usepackage{array}
\usepackage{booktabs}
\usepackage[hidelinks]{hyperref}

\definecolor{accent}{HTML}{1F4E79}
\definecolor{ink}{HTML}{222222}
\definecolor{muted}{HTML}{6B6B6B}
\pagestyle{empty}
\setlength{\parindent}{0pt}
\newcolumntype{L}{>{\raggedright\arraybackslash}X}
% \fieldname{Bill to}: a small-caps label.
\newcommand{\fieldname}[1]{\textcolor{muted}{\footnotesize\scshape\MakeLowercase{#1}}}

\begin{document}
\color{ink}
{% assign d = currency | currency_decimals %}{% assign subtotal = 0 %}{% for item in items %}{% assign price = item.unit_price | round: 4 %}{% assign line = item.quantity | times: price | round: d %}{% assign subtotal = subtotal | plus: line %}{% endfor %}{% assign tax = subtotal | times: tax_rate | divided_by: 100.0 | round: d %}{% assign total = subtotal | plus: tax %}
\begin{tabularx}{\linewidth}{@{}L r@{}}
  {\Large\bfseries\mbox{}{{ seller.name }}} & \textcolor{accent}{\fontsize{24}{26}\selectfont\scshape invoice} \\[6pt]
  \parbox[t]{0.55\linewidth}{\raggedright\small\mbox{}{{ seller.address }}\par\mbox{}{{ seller.email }}} &
  {\small\begin{tabular}[t]{@{}r@{\hspace{8pt}}l@{}}
    \textcolor{muted}{Invoice no.} & {{ invoice_number }} \\
    \textcolor{muted}{Date} & {{ issue_date }} \\
    \textcolor{muted}{Due} & {{ due_date }} \\
    \textcolor{muted}{Terms} & {{ payment_terms }}
  \end{tabular}}
\end{tabularx}

\vspace{10pt}
\textcolor{accent}{\rule{\linewidth}{0.8pt}}
\vspace{6pt}

\fieldname{Bill to}\par\vspace{2pt}
{\bfseries\mbox{}{{ client.name }}}\par
\parbox[t]{0.55\linewidth}{\raggedright\mbox{}{{ client.address }}}

\vspace{18pt}
% A longtable, so a long list breaks across pages with the header repeated.
% Every column has a fixed width: one compile pass gives aligned columns.
{\renewcommand{\arraystretch}{1.45}\setlength{\LTpre}{0pt}\setlength{\LTpost}{0pt}\setlength{\LTleft}{0pt}\setlength{\LTright}{0pt}\arrayrulecolor{accent}
\begin{longtable}{@{}>{\raggedright\arraybackslash}p{\dimexpr\linewidth-6\tabcolsep-12mm-26mm-26mm\relax} >{\raggedleft\arraybackslash}p{12mm} >{\raggedleft\arraybackslash}p{26mm} >{\raggedleft\arraybackslash}p{26mm}@{}}
  \toprule
  \textbf{Description} & \textbf{Qty} & \textbf{Unit price} & \textbf{Amount} \\
  \midrule
\endhead
  \bottomrule
\endfoot
  \bottomrule
\endlastfoot
{% for item in items %}{% assign price = item.unit_price | round: 4 %}{% assign line = item.quantity | times: price | round: d %}  {{ item.description }} & {{ item.quantity | round: 2 }} & {{ price | currency: currency, 4 }} & {{ line | currency: currency }} \\
{% endfor %}\end{longtable}}

\vspace{8pt}
\hfill{\renewcommand{\arraystretch}{1.35}\begin{tabular}{@{}l r@{}}
  \textcolor{muted}{Subtotal} & {{ subtotal | currency: currency }} \\
  \textcolor{muted}{Tax ({{ tax_rate }}\%)} & {{ tax | currency: currency }} \\
  \arrayrulecolor{accent}\specialrule{0.8pt}{2pt}{2pt}
  \textbf{Total due} & \textcolor{accent}{\large\bfseries {{ total | currency: currency }}} \\
\end{tabular}}

\vfill
\textcolor{muted}{\rule{\linewidth}{0.4pt}}\par\vspace{6pt}
\begin{tabularx}{\linewidth}{@{}L@{\hspace{12mm}}L@{}}
  \fieldname{Payment} & \fieldname{Notes} \\
  {\small {{ payment_details }}} & {\small {{ notes }}}
\end{tabularx}
\end{document}
